\section{Results}
\label{sec:results}

\begin{table*}[t]
  \centering
  \caption{Main results over \NumScoredEdits{} scored edits (\NumFacesScored{}
  faces, \NumRunsScored{} runs). Medians per cell; identity shows min--max in
  parentheses. Loc.\ is the localization ratio; $\Delta E_{\mathrm{tgt}}$ the mean
  CIELAB change inside the target region. Chained rows are single-face
  order-effect probes, not population estimates.}
  \label{tab:main}
  % Generated by scripts/aggregate.py - do not edit by hand.
\begin{tabular}{lllrlllrr}
\toprule
Procedure & Model & Control & $n$ & Identity$\uparrow$ & Loc.$\uparrow$ & $\Delta E_{\mathrm{tgt}}$ & Lat.\,(s) & Cost\,(\$) \\
\midrule
Deep plane facelift & GPT Image 2 & Masked composite & 5 & 0.901 {\scriptsize(0.855--0.964)} & 0.997 & 4.2 & 63 & 0.220 \\
Deep plane facelift & GPT Image 2 & Prompt only & 4 & 0.883 {\scriptsize(0.855--0.924)} & 0.553 & 3.9 & 64 & 0.220 \\
Chained: facelift $\to$ rhino. & GPT Image 2 & Masked composite & 1 & 0.841 {\scriptsize(0.841--0.841)} & 0.984 & 9.5 & -- & -- \\
Rhinoplasty & GPT Image 2 & Masked composite & 15 & 0.837 {\scriptsize(0.643--0.951)} & 0.998 & 6.1 & 70 & 0.220 \\
Rhinoplasty & GPT Image 2 (low) & Masked composite & 8 & 0.887 {\scriptsize(0.786--0.976)} & 0.997 & 4.3 & 38 & 0.015 \\
Rhinoplasty & Qwen-Image-Edit (inpaint) & Masked inpaint & 8 & 0.973 {\scriptsize(0.873--0.990)} & 0.530 & 1.9 & 5 & 0.030 \\
Rhinoplasty & Seedream 5.0 & Masked composite & 10 & 0.926 {\scriptsize(0.786--0.962)} & 0.999 & 10.0 & 46 & 0.035 \\
Chained: rhino.\ $\to$ facelift & GPT Image 2 & Masked composite & 1 & 0.848 {\scriptsize(0.848--0.848)} & 0.981 & 10.1 & -- & -- \\
\bottomrule
\end{tabular}

\end{table*}

\begin{figure}[t]
  \centering
  \includegraphics[width=\columnwidth]{scatter_identity_localization}
  \caption{Every scored edit, placed by the two failure modes. Masked compositing
  pins localization near 1 across models; prompt-only edits and the hosted
  inpainter leak roughly half their change off-target. All edits clear the
  identity floor.}
  \label{fig:scatter}
\end{figure}

Table~\ref{tab:main} and Fig.~\ref{fig:scatter} summarize the \NumScoredEdits{}
scored edits. We report medians and ranges throughout and make no statistical
significance claims at this sample size.

\subsection{Localization: the ladder works}

Masked compositing drives the localization ratio to a median of
\MedLocComposite{} (range \LocCompositeMin--\LocCompositeMax) for all six
editors. The same models under prompt-only control reach only
\MedLocPromptOnly{} (\LocPromptOnlyMin--\LocPromptOnlyMax): roughly half of the
measured change lands outside the surgical region, which is the beauty-filter
signature the metric was built to expose. The off-target numbers make the contrast
concrete---a median $\Delta E_{\mathrm{off}}$ of \MedOffPromptOnly{} for
prompt-only edits against \MedOffComposite{} for composited ones---while the
on-target change stays substantial (median $\Delta E_{\mathrm{tgt}}$
\MedTgtComposite{} composited vs.\ \MedTgtPromptOnly{} prompt-only). Prompt
guardrails alone, however carefully worded, did not confine any model's edit;
moving the confinement outside the model did, uniformly.

\subsection{Identity}

With six editors in the matrix the identity floor becomes a discriminating bar
rather than a formality. \NumBelowFloor{} of \NumScoredEdits{} edits fall below
$0.6$, and \NumBelowFloorFlux{} of them come from FLUX.2~[pro], which sits at one
end of an aggressiveness spectrum: it produces the largest anatomical changes in
the study (median prompt-only $\Delta E_{\mathrm{tgt}}$ \MedTgtFluxPromptOnly)
and pays for them in identity (prompt-only median
\MedIdentityFluxPromptOnly{}, single edits as low as \IdentityMinAll{}---at that
score the output is effectively a different person). The Nano Banana models
occupy the opposite end, holding identity at \MedIdentityNanoBanana{} while
making the gentlest edits (median $\Delta E_{\mathrm{tgt}}$ \MedTgtNanoBanana).
GPT Image~2 and Seedream~5.0 fall between. Following
Sec.~\ref{sec:protocol}, the model-level number is the prompt-only median of
\MedIdentityPromptOnly{}; composited edits score \MedIdentityComposite{}, but
part of that is copied pixels, so we do not read the difference as a model
improvement. Fig.~\ref{fig:boxes} breaks identity out by model; the ordering
tracks each model's editing aggressiveness rather than output quality.

\begin{figure*}[t]
  \centering
  \includegraphics[width=0.92\textwidth]{box_localization_identity}
  \caption{(a) Localization by control rung, models pooled. (b) Identity by
  model, controls pooled; the dashed line is the acceptance floor. FLUX.2's low,
  wide distribution in (b) is the aggressiveness trade of
  Sec.~\ref{sec:results}-B; the inpainter's spread reflects its per-procedure
  split (Sec.~\ref{sec:inpaint}).}
  \label{fig:boxes}
\end{figure*}

\subsection{Hosted inpainting under-edits}
\label{sec:inpaint}

The masked-inpaint rung was expected to dominate: the endpoint receives the mask
explicitly and should confine the edit for free. Instead it fails differently on
each procedure. On rhinoplasty it shows the did-nothing trap: identity of
\InpaintRhinoIdentity{} with a target change of only \InpaintRhinoTgt{}---barely
above its off-target change of \InpaintRhinoOff{}, for a localization of
\InpaintRhinoLoc{}. The nose edit is timid while the returned image differs
slightly everywhere, consistent with a full-image re-encode by the endpoint. On
the facelift the same endpoint does edit (target change \InpaintFaceliftTgt{},
localization \InpaintFaceliftLoc{}) but drops identity to
\InpaintFaceliftIdentity{}, the lowest median of any rung on that procedure, and
one of its eight requests was rejected by the provider's content checker. Either
way the conclusion holds: at the tested endpoint, sending the mask to the model
was worse than ignoring the model's mask support and compositing on the client.

\subsection{Ground-truth sanity check}

Against real post-operative photographs the edits score a median cosine of
\MedGtCosine{} (range \GtCosineMin--\GtCosineMax; facelift
\MedGtCosineFacelift{}, rhinoplasty \MedGtCosineRhino), shown in
Fig.~\ref{fig:gt}. The values sit where same-person photographs taken months apart
under different pose, makeup, and hair typically land, which is all this check can
certify: the previews are not drifting toward a different person than the one the
surgery actually produced.

\begin{figure}[t]
  \centering
  \includegraphics[width=\columnwidth]{strip_gt_cosine}
  \caption{ArcFace cosine between each edit and the face's real post-operative
  photograph (bars: medians). A soft sanity check; session differences cap the
  achievable score.}
  \label{fig:gt}
\end{figure}

\subsection{Chained procedures}

Chaining rhinoplasty and the facelift on the same face works in both orders:
identity \ChainedIdentityLow--\ChainedIdentityHigh{} and localization
\ChainedLocLow--\ChainedLocHigh{} for the two orderings ($n=\NumChained$), with no
visible interference between the nose and jaw--neck regions. Because the regions
are disjoint, each composited stage protects the other stage's canvas; order
appears to be a free parameter, though two probes on one face only establish
feasibility.

\subsection{Qualitative results and failure modes}

Fig.~\ref{fig:teaser} and Fig.~\ref{fig:grid} show representative edits. The
facelift examples display the intended tissue repositioning---jawline definition,
reduced jowls, a cleaner cervicomental angle---while skin texture, lighting, and
background survive untouched in the composited column. Three failure modes
recurred. The ghost seam: before pose gating, profile views put the frontal-mesh
mask on the cheek and composited a translucent copy of the nose into it; gating
on yaw removed the artifact at the cost of leaving profile inputs uncomposited.
The person swap: FLUX.2's strongest edits replace enough of the face to break
identity, and compositing can only partially repair this because the damage is
inside the mask. The detector abstention: on one extreme-profile input the identity model finds
no face at all, so that face's twelve cells contribute ground-truth cosines but
no identity scores---counted, not hidden, in Table~\ref{tab:main}.

\begin{figure*}[t]
  \centering
  \includegraphics[width=0.88\textwidth]{qualitative_grid.png}
  \caption{Qualitative results (top to bottom: two deep-plane-facelift faces, one
  frontal rhinoplasty face). Prompt-only edits alter the whole image; composited
  edits are pixel-identical to the input outside the region.}
  \label{fig:grid}
\end{figure*}
